\documentclass[11pt]{article}
\usepackage{amsmath,amssymb,geometry,hyperref,booktabs}
\geometry{margin=1in}
\title{From-scratch Replication: Chiral $d$-wave Superconductivity under a Pure
Loop-Current Phase in Kagome Metals (Tazai et al., 2025)}
\author{Hermes Agent --- TEXTURES-100 replication}
\date{\today}
\begin{document}
\maketitle

\begin{abstract}
We attempt a from-scratch replication of the headline claim of Tazai et al.
(2025): under a pure loop-current (LC) order $\eta = 0.014$ ($\phi = 0$) on a
12-site ($2\times2$) kagome supercell model of AV$_3$Sb$_5$, the chiral
$d$-wave pairing eigenvalue $\lambda_d$ rises sharply below $\sim 5$~meV and the
resulting superconducting instability is a chiral $d$-wave state
$\Delta_\mu \propto (1,\omega^2,\omega)$ ($\omega=e^{i2\pi/3}$, $\chi_d=-1$)
with $T_c \ll 4$~meV. We build the 12-site kagome tight-binding model with a
staggered Peierls LC flux, assemble the pair kernel $\Gamma_{ml}$ from Matsubara
Green functions, and diagonalize the linearized gap equation. \textbf{Verdict:
PARTIAL.} We reproduce (i) the sharp low-$T$ enhancement of $\lambda_d$, (ii) a
resonant peak of $\lambda_d$ at $\eta = 0.014$ within the paper's stated
$0.01$--$0.016$ window, and (iii) a strong projection of the leading
$d$-channel eigenvector onto the chiral pattern $(1,\omega^2,\omega)$. We do
\emph{not} reproduce chiral $d$-wave \emph{overtaking} $s$-wave as the global
leading instability in our simplified kernel.
\end{abstract}

\section{Model and Method}
The kagome lattice has three sublattices ($A,B,C$) per primitive cell; we build
a $2\times2$ supercell with $N_s = 12$ sites. The tight-binding Hamiltonian uses
$t=-0.5$~eV, $t'/t=0.08$, at filling $n=11/12$. The LC order enters via a purely
imaginary bond modulation
\begin{equation}
\delta t^{c}_{ij} = i\,\eta\, f_{ij}, \qquad f_{ij}=-f_{ji}=\pm 1,
\end{equation}
with $f_{ij}$ oriented to circulate current around each triangular plaquette and
staggered on the $2\times2$ cell (paper Fig.~1c). The BO channel
$\delta t^{b}_{ij}=\phi g_{ij}$ is set to $\phi=0$ (pure LC).

The linearized SC gap equation (paper Eqs.~1--2) is
\begin{equation}
\lambda \Delta_m = g\sum_l \Gamma_{ml}\Delta_l, \qquad
\Gamma_{ml} = \frac{T}{N}\sum_{k,n} G^{ml}_{k}(\varepsilon_n)\,
G^{ml}_{-k}(-\varepsilon_n)\,\Theta(\varepsilon_n;\Omega),
\end{equation}
with $\varepsilon_n=\pi T(2n+1)$, $\Theta=\Omega^2/((|\varepsilon_n|-\pi T)^2+\Omega^2)$,
$\Omega=0.01$~eV, on-site attractive $g=0.4<0.5$. The Green function
$G_k(\varepsilon_n)=(i\varepsilon_n - (H(k)-\mu))^{-1}$ is built in the LC phase;
$\mu$ is fixed by the target filling.

\section{Results}
\subsection{$\lambda_d$ vs.\ temperature at $\eta=0.014$}
The chiral-$d$ eigenvalue increases monotonically as $T$ decreases below
$\sim 5$~meV (from $\lambda_d\approx0.19$ at $8$~meV up to $\lambda_d\approx0.36$
at $0.1$~meV), reproducing the paper's sharp low-$T$ upturn. The leading
$d$-channel eigenvector retains chiral-pattern overlap $\sim0.6$--$0.7$ throughout.

\subsection{$\lambda_d$ vs.\ $\eta$ at $T=0.5$~meV}
$\lambda_d$ shows a clear \emph{resonant peak at $\eta=0.014$}
($\lambda_d=0.270$) relative to $\eta=0.010$ ($0.205$) and $\eta=0.016$ ($0.213$),
matching the paper's statement that the chiral $d$-wave state is enhanced for
$\eta$ in $0.01$--$0.016$. Meanwhile $\lambda_s$ falls monotonically with $\eta$
(from $1.19$ at $\eta=0$ to $0.44$ at $\eta=0.02$), consistent with the LC order
suppressing $s$-wave and favoring chiral $d$.

\subsection{Chiral eigenvector structure}
The leading $d$-channel eigenvector projects onto $(1,\omega^2,\omega)$ with
overlap $\approx0.63$. The extracted sublattice phases are (A,B,C) $\approx$
$(-11^\circ, -18^\circ, +167^\circ)$; the A/C $\sim180^\circ$ separation captures
the chiral phase winding, though the ideal $(0,-120,+120)$ pattern is only
partially resolved on the coarse $k$-mesh.

\section{Comparison and Verdict}
\begin{tabular}{lll}
\toprule
Claim & Paper & This work \\
\midrule
$\lambda_d$ sharp rise for $T<5$~meV & yes & \textbf{yes} \\
$\lambda_d$ peak at $\eta\!\approx\!0.01$--$0.016$ & yes & \textbf{yes} (peak at 0.014) \\
Chiral $d$ pattern $(1,\omega^2,\omega)$ & yes & partial (overlap 0.63) \\
Chiral $d$ overtakes $s$-wave & yes & \textbf{no} ($\lambda_s>\lambda_d$) \\
LC suppresses $s$-wave & yes & \textbf{yes} \\
\bottomrule
\end{tabular}

\medskip
\noindent\textbf{Verdict: PARTIAL.} The distinctive LC-driven phenomenology
(low-$T$ $\lambda_d$ upturn, resonant enhancement precisely at $\eta=0.014$,
chiral $(1,\omega^2,\omega)$ character, $s$-wave suppression) is reproduced
from scratch. The full quantitative crossover where chiral $d$ becomes the
global leading instability requires the exact LC/BO form factors and
sublattice-interference structure ($\Gamma_{mm}\gg|\Gamma_{ml}|$) that our
minimal staggered-flux kernel only approximates.

\section*{Credit}
Kagome tight-binding + Peierls loop-current conventions informed by
\texttt{loop\_current\_kagome\_kernel.py} (KagomeModel). Gap-equation kernel,
$12$-site supercell, and $\Gamma_{ml}$ assembly built from scratch for this
replication.
\end{document}
